\section{Related Work}\label{sec:related}
\answers{where RQ1--RQ7 sit relative to prior work}

\paragraph*{Prompt-injection benchmarks for agents} Greshake et~al.\ introduced indirect prompt injection against LLM-integrated applications~\cite{greshake2023}, and Liu et~al.\ formalised attacks and defences~\cite{liu2024formalizing}. InjecAgent~\cite{injecagent} and AgentDojo~\cite{agentdojo} benchmark tool-using agents against injected tool outputs, Agent Security Bench~\cite{asb2025} broadens the attack and defence set, ToolEmu~\cite{toolemu} emulates tools to find risky agent behaviour, and RedCode~\cite{redcode2024} benchmarks risky code execution by code agents. These fix the agent framework and vary model and attack; our bench fixes model and payload and varies real harness binaries and modes.

\paragraph*{Design-level defences} CaMeL~\cite{camel2025} and the design patterns of Beurer-Kellner et~al.~\cite{designpatterns2025} constrain what an agent can do after reading untrusted data. System-level isolation (IsolateGPT~\cite{isolategpt2025}) and programmable privilege control (Progent~\cite{progent2025}) move enforcement into the runtime, much as our rubric asks harnesses to do. Willison's ``lethal trifecta''~\cite{willison2025trifecta} explains why egress (D03) and secret handling (D07) matter even when injection cannot be prevented.

\paragraph*{MCP security} Hou et~al.\ survey the MCP landscape and its threats~\cite{hou2025mcp}. Radosevich and Halloran show that MCP-connected models can be driven to harmful actions~\cite{radosevich2025mcp}. Hasan et~al.\ study the security and maintainability of MCP servers at scale~\cite{hasan2025mcp}, and MCPTox benchmarks tool poisoning on real servers~\cite{mcptox2026}. Invariant Labs described tool-poisoning attacks and the GitHub MCP data leak~\cite{invariant2025toolpoisoning,invariant2025github}, and Trail of Bits showed that servers can attack before any tool is used~\cite{tob2025mcpline}. We add the harness side: how servers are declared, pinned and approved, and how often risky declarations appear in public repositories.

\paragraph*{Practitioner research on coding agents} Much of what is known comes from practitioners: Pillar Security's rules-file backdoor~\cite{pillar2025rulesfile}, Tracebit's Gemini CLI allowlist bypass~\cite{tracebit2025gemini}, Johann Rehberger's \emph{Embrace The Red} series, including the Copilot auto-approve RCE, the Month of AI Bugs and the Antigravity analysis~\cite{etr2025copilot,etr2025monthofaibugs,etr2025antigravity}, Trail of Bits on prompt injection to RCE in agents~\cite{tob2025agentrce}, and Check Point Research and Cymulate on Cursor, Codex CLI and Claude Code~\cite{checkpoint2025mcpoison,checkpoint2025codex,checkpoint2026claude,cymulate2025inverse}. Our survey normalises these reports into one taxonomy, root-cause scheme and counting rule.

\paragraph*{CI and supply chain} Studies of GitHub workflow security~\cite{koishybayev2022ci,argus2023} underpin our CI rules and E2 workflow stratum, and open-source supply-chain attack taxonomies~\cite{ohm2020npm} frame E4.

\paragraph*{Frameworks} The OWASP Top~10 for LLM Applications~\cite{owaspllm2025}, the OWASP Agentic AI threats and mitigations guide~\cite{owaspagenticthreats2025}, the OWASP Top~10 for Agentic Applications~\cite{owaspagentic2026} and the OWASP MCP Top~10~\cite{owaspmcp} catalogue risks at the application level, and MITRE ATLAS~\cite{atlas} catalogues adversary techniques. We map our taxonomy onto all of them and add a harness-specific, checkable control layer.

\paragraph*{Author's prior tooling} \tool{} reuses parsing ideas from the author's \texttt{upfront}, offers \texttt{rulewarden} (deterministic instruction-file rule enforcement) as an integration point, and uses \texttt{airmkit}'s control vocabulary for the MAS mapping.
